In \autoref{ch:boolean-model}, a first experiment on TF activities and PD-L1 expression
was conducted using a method designed for binary data.
One drawback of this method is that, after binarization,
information on which binary features belong to which TF is lost.
This information could be used in the model, for instance to speed up computation.

The TF activity scores are continuous.
Instead of pre-binarizing via thresholding and employing the boolean-featured model,
we adjust the boolean model to consider threshold clauses on continuous features.
This can not only speed up solving,
it sets the stage for future extensions that attempt to exploit the continuous nature of the TF data.
This model, which we call the threshold model, is a simple adjustment of the boolean model and is our contribution.

The structure of this chapter follows closely that of \autoref{ch:boolean-model}.
A reader who is not interested in the mathematics and implementation of the method but rather the next experiments and their results can
skip to \autoref{sec:threshold-methodology}.

\section{Model description}\label{sec:threshold-model-description}
Recall the boolean model as described in \autoref{sec:boolean-model-description}.
In this setting, feature values are continuous, so $\bm{X}_j \in \R^\J, \; j \in \J$.
A clause represents a lower-bound, and consists of a feature and a threshold.
An upper-bound is equivalent to a lower-bound on the negated feature, i.e.
\[X_{\cdot j} \leq \alpha \iff -X_{\cdot j} \geq -\alpha.\]
Hence, without loss of generality, only lower-bounds are considered.
We denote the set of possible clauses by $\cL \subseteq \J \times \mathbb{R}.$
The candidate rule set remains $\K = 2^\L$, the power set of $\cL$.
In this setting,
\begin{equation}
    \K_i = \left\{k \in \K : \underset{(j,\alpha) \in k}{\bigwedge} X_{ij} \geq \alpha\right\}, \tquad i \in \I. \label{eq:threshold_K_i}
\end{equation}
The master problem formulation \eqref{eq:BM} remains unchanged.

While a threshold can theoretically be any value,
it suffices to consider only thresholds $\{X_{ij}\}_{i \in \I}$ for any feature $j$.
Any other threshold yields a clause that is equivalent, at least in training, to one already considered.
For example, if \[\J = \{a\},\; \I=\{1,2\} \text{ and } X_{1 a} = 1, \; X_{2 a} = 2,\] the only relevant lower-bounds are
\[X_{\cdot a} \geq 1 \quad \text{and} \quad X_{\cdot a} \geq 2.\]
Hence, we consider
\begin{equation}
    \cL = \left\{\left(j, X_{ij}\right) : j \in \J, \; i \in \I\right\}.\label{eq:threshold_L_full}
\end{equation}
The pricing problem IP is similar to \eqref{eq:BP}.
The difference is that we now speak of $z_l, \; l \in \L$, and we define $S_i$ as
\begin{equation}\label{eq:S_i_threshold}
    S_i = \{(j,\alpha) \in \cL : X_{ij} < \alpha\}, \tquad i \in \I.
\end{equation}
Hence, $S_i$ the set of clauses that cut off point $i$.
The IP is given by
\begin{subequations}
    \label{eq:TP}
    \begin{align}
        \min \quad & \sum_{i \in \cN} \delta_i - \sum_{i \in \cP} \mu_i \delta_i +
        \lambda c(\bm{z}) \label{eq:TP-obj} \\
        \st \quad & D\delta_i + \sum_{l \in S_i} z_l \leq D, & i \in \I, \label{eq:TP-con1}\\
        & \delta_i +  \sum_{l \in S_i} z_l \geq  1, & i \in \I,  \label{eq:TP-con2}\\
        & \sum_{l \in \L} z_l \leq D, \quad \label{eq:TP-complex}\\
        & \mathrlap{\delta\in \{0,1\}^\I,~z\in\{0,1\}^{\cL}.}\label{eq:TP-vars}
    \end{align}
\end{subequations}

We impose another constraint in this setting: the use of at most one clause per feature $j \in \J$.
This constraint can be written as
\begin{equation}
    \sum_{(j',\alpha) \in \cL \,:\, j' = j} z_{(j', \alpha)} \leq 1, \tquad \quad j \in \J. \label{eq:TP_one-per-feature}
\end{equation}
If negated feature values are added as separate features to model upper-bounds,
then \eqref{eq:TP_one-per-feature} does not prevent a lower-bound and an upper-bound on the same feature from being used in one rule.
To achieve this, a separate constraint family may be used.

Note that $|\L| = |\I||\J|.$ In our case, $|\L| \approx 800 \cdot 20 \approx 16 000.$
This is rather large, causing $|\K|$ to explode.
An alternative is to choose
\begin{equation}\label{eq:threshold_L_quantile}
    \cL = \left\{ \left(j, Q_{\bm{X}_j}(q)\right) : j \in \J, \; q \in \bm{q} \right\}.
\end{equation}
where $\bm{q}$ is a vector of predefined quantile levels.
In particular, given a number of quantiles $\nq$, one can consider the quantile level vector
\[ \bm{q} = \left(\frac{\nu}{\nq + 1}\right)_{\nu = 1,\dots,\nq},\]
as chosen in \autoref{sec:boolean-methodology}.
This approach is equivalent to pre-binarizing and using the boolean-featured model.

To simplify interpretability of and control over the model, we define rule complexity as
\[c_k = 1, \tquad k \in \K.\]
This implies that $C$ is the maximum cardinality of a ruleset.

\section{Implementation}\label{sec:threshold-implementation}
The threshold model as described in \autoref{sec:threshold-model-description}
is a special case of the ramp loss model.
We refer the reader to \autoref{ch:ramp_model} for details.

\section{Methodology}\label{sec:threshold-methodology}
First, we focus on simple rulesets for low PD-L1 expression thresholds.
High values of $Q_y$ led to very little learning as discussed in \autoref{sec:boolean-results}.
Therefore, we restrict ourselves to $Q_y = 0.5, 0.6.$
We are interested in the performance of very simple rulesets and what they look like.
Therefore, we study
\[C = 1,2 \quad \text{and} \quad D = 1,2,3,\]
training on Broad in full.
We fix $\nq = 9$ and consider both lower-bounds and upper-bounds.
The use of the exact pricer has proven to be time-consuming, especially for higher $\nq$.
We are interested in high $\nq$ to rule out the effect of course discretization on ruleset performance.
To avoid wasting time and compute (that could otherwise be used by other users of the HPC),
we use only the heuristic pricer with beam widths $(1000,1000)$ for $D=2$ and $(1000,1000,1000)$ for $D=3.$
While this removes optimality guarantees like those discussed in \autoref{sec:boolean_model_appendix},
rulesets still provide a proof that a certain performance can be attained.
Furthermore, optimal solutions for small $C$ and small $D$ will be discussed in \autoref{sec:threshold_greedy}.

The total time limit is set to 3 hours, which is much more than needed for this setup.

Next, we study the effect of increasing $C$ and $D$ on performance.
Again, Broad is used in full for training.
The settings for this experiment are analogous to the experiment above.

Finally, we study the generalization of simple rulesets for $Q_y=0.5, 0.6.$
We apply 10-fold stratified cross-validation on Broad.
Then we compare train and test accuracy.
We set $D = 2$ and study $C=1,2,3$.
We also study the consistency of these solutions by counting how often each ruleset occurs among splits.
In this count, we disregard the threshold value and direction (lower-bound or upper-bound) of each clause,
but identify a ruleset only by which TFs appear together in each rule.
Ideally, the ruleset found is consistent, as it is meant to be a general rule governing PD-L1 expression.

\section{Results}\label{sec:threshold-results}
The results for $C=1$ are summarized in \autoref{tab:threshold_interpretable_results_C_1}.
We report accuracy (Acc.), precision (Prec.) and recall (Rec.).
A threshold is shown as both an absolute value (i.e. $-0.20$) and a quantile level
with respect to the observed activity scores of this TF in the training dataset (i.e. $Q(0.50)$).
With merely one rule, a remarkable accuracy can be achieved.
Increasing $D$ consistently yields the same rule, plus a new clause.
Furthermore, the rulesets found for $Q_y=0.5$ and $Q_y=0.6$ are distinct.
There is a performance jump going from $D=1$ to $D=2$, less so going from $D=2$ to $D=3$.
This suggests considering a compound rule, rather than simply a clause, has some effect, but
adding more clauses in the compound rule does not appear to have much effect.

The results for $C=2$ are summarized in \autoref{tab:threshold_interpretable_results_C_2}.
By allowing for two rules, accuracy only increases marginally.
Increasing $D$ from $2$ to $3$ again yields practically no performance gain.
Notably, increasing $D$ now yields qualitatively different rulesets.
Another observation is that for $Q_y=0.5$, $C=2$, $D=2$, both rules involve the same TF\@.

\begin{table}[h]
    \centering
    \bwtable{gen/PDL1ex2/PDL1ex2_tab_interpretable_mw_0_C_1}
    \caption{Results for the full Broad dataset for $C=1$ and small rules ($D=1,2,3$).}
    \label{tab:threshold_interpretable_results_C_1}
\end{table}

\begin{table}[h]
    \bwtable{gen/PDL1ex2/PDL1ex2_tab_interpretable_mw_0_C_2}
    \caption{Results for the full Broad dataset for $C=2$ and small rules ($D=1,2,3$).}
    \label{tab:threshold_interpretable_results_C_2}
\end{table}

The results for the second experiment are shown in \autoref{fig:threshold_accuracy_mw_0.svg}.
First observe that setting $C=5, \; D=3$ in the threshold model
is comparable to setting $C=20, \; D=3$ in the boolean model.
Therefore, compared to the experiment described in \autoref{sec:boolean-methodology},
this experiment considers lower ruleset complexity.
As seen in previous tables, low values of $C$ yield reasonable performance.
While accuracy increases with $C$, the increase is marginal considering the increase in perceived ruleset complexity.
To clarify, a table like \autoref{tab:threshold_interpretable_results_C_1} for $C=5, D=3$ would span several pages.
Increasing $D$ has little effect on performance and even $D=1$ shows reasonable performance for low $Q_y$.
However, as seen in the previous experiment, there is some performance jump going from $D=1$ to $D=2$.

The cross-validation results for $Q_y=0.5, 0.6$ and $D=2$ are shown in
\autoref{fig:threshold_train_test_accuracy_mw_0.svg}.
If $C=1$, then train and test accuracy are close.
Increasing $C$ marginally improves training accuracy, but does not significantly improve test accuracy.
Hence, it appears that increasing $C$ causes overfitting rather than learning.

At this point, it appears that the most effective rulesets are the absolute simplest ones:
those with one rule and only one or two TF\@.
It may be that more complicated rulesets simply do not perform better on this dataset.
It may also be that the model does not behave as intended.
Hence, we study an alternative approach with an independent Python implementation in \autoref{sec:threshold_greedy}.

Finally, the ruleset frequency counts are given in
\autoref{tab:threshold_cross_interpretable_mw_0_C_1}.
This reveals that rulesets are sometimes (nearly) consistent, but not always.
A possible explanation for the inconsistency is that some TF are known to bind, for instance RELA and NFKB1.
It is therefore possible that the complex claimed to be promoting PD-L1 expression is a complex of
RELA, NFKB1 and STAT2.
However, NFKB1 only appears once in the entire table.
More generally, the high cross-correlation between TF activity scores as observed in \autoref{fig:corr_activity}
suggests that TF could be effectively interchangeable in the dataset.
Whether this is a real-life phenomenon or an artifact of the TF estimation technique is unclear.

\begin{figure}[h]
    \centering
    \includesvg[width=0.7\textwidth]{gen/PDL1ex2/PDL1ex2_fig_accuracy_mw_0.svg}
    \caption{Accuracy of results for the full Broad dataset for $C=1,2,3,4,5$, $D=1,2,3,4$ and $Q_y=0.5,0.6,0.7,0.8,0.9.$}
    \label{fig:threshold_accuracy_mw_0.svg}
\end{figure}

\begin{figure}[h]
    \centering
    \includesvg[width=0.8\textwidth]{gen/PDL1ex2/fig_train_test_accuracy_mw_0.svg}
    \caption{Confidence intervals ($95\%)$ for train and test accuracy of results under 10-fold stratified cross-validation on the Broad dataset, $C=1,2,3.$}
    \label{fig:threshold_train_test_accuracy_mw_0.svg}
\end{figure}

\begin{table}[h]
    \centering
    \begin{tabular}{llllr}
\toprule
 &  &  &  & \textbf{Frequency} \\
$\mathbf{Q_y}$ & $\mathbf{C}$ & $\mathbf{D}$ & \textbf{TF combination} &  \\
\midrule
\multirow[t]{6}{*}{0.5} & \multirow[t]{6}{*}{1} & \multirow[t]{3}{*}{2} & {RELA $\wedge$ STAT2} & 6 \\
 &  &  & {FOS $\wedge$ STAT2} & 3 \\
 &  &  & {NFKB1 $\wedge$ STAT2} & 1 \\
\cline{3-5}
 &  & \multirow[t]{3}{*}{3} & {FOS $\wedge$ IRF1 $\wedge$ STAT2} & 5 \\
 &  &  & {FOS $\wedge$ RELA $\wedge$ STAT2} & 3 \\
 &  &  & {IRF1 $\wedge$ RELA $\wedge$ STAT2} & 2 \\
\cline{1-5} \cline{2-5} \cline{3-5}
\multirow[t]{3}{*}{0.6} & \multirow[t]{3}{*}{1} & 2 & {FOSL1 $\wedge$ STAT2} & 10 \\
\cline{3-5}
 &  & \multirow[t]{2}{*}{3} & {FOS $\wedge$ FOSL1 $\wedge$ STAT2} & 9 \\
 &  &  & {FOSL1 $\wedge$ STAT2 $\wedge$ STAT6} & 1 \\
\cline{1-5} \cline{2-5} \cline{3-5}
\bottomrule
\end{tabular}

    \caption{Ruleset frequency count for the results of 10-fold stratified cross validation on the Broad dataset.
    In the frequency count, the threshold values and directions have been disregarded.
    Therefore, this table shows how often distinct TF combinations appear when only one rule is to be deduced from a cross-validation split of the Broad dataset}
    \label{tab:threshold_cross_interpretable_mw_0_C_1}
\end{table}
%
%\begin{table}[h]
%    \centering
%    \begin{tabular}{llllr}
\toprule
 &  &  &  & \textbf{Frequency} \\
$\mathbf{Q_y}$ & $\mathbf{C}$ & $\mathbf{D}$ & \textbf{TF combination} &  \\
\midrule
\multirow[t]{4}{*}{0.5} & \multirow[t]{4}{*}{1} & \multirow[t]{2}{*}{2} & {RELA $\wedge$ STAT2} & 7 \\
 &  &  & {FOS $\wedge$ STAT2} & 3 \\
\cline{3-5}
 &  & \multirow[t]{2}{*}{3} & {FOS $\wedge$ RELA $\wedge$ STAT2} & 9 \\
 &  &  & {FOS $\wedge$ IRF1 $\wedge$ STAT2} & 1 \\
\cline{1-5} \cline{2-5} \cline{3-5}
\multirow[t]{5}{*}{0.6} & \multirow[t]{5}{*}{1} & \multirow[t]{2}{*}{2} & {FOSL1 $\wedge$ STAT2} & 9 \\
 &  &  & {RELA $\wedge$ STAT2} & 1 \\
\cline{3-5}
 &  & \multirow[t]{3}{*}{3} & {FOS $\wedge$ FOSL1 $\wedge$ STAT2} & 8 \\
 &  &  & {FOSL1 $\wedge$ STAT2 $\wedge$ STAT6} & 1 \\
 &  &  & {FOSL1 $\wedge$ RELA $\wedge$ STAT2} & 1 \\
\cline{1-5} \cline{2-5} \cline{3-5}
\bottomrule
\end{tabular}

%    \label{tab:threshold_cross_interpretable_mw_0.2_C_1}
%\end{table}
%
%\begin{table}[h]
%    \centering
%    \begin{tabular}{llllr}
\toprule
 &  &  &  & \textbf{Frequency} \\
$\mathbf{Q_y}$ & $\mathbf{C}$ & $\mathbf{D}$ & \textbf{TF combination} &  \\
\midrule
\multirow[t]{3}{*}{0.5} & \multirow[t]{3}{*}{1} & 2 & {RELA $\wedge$ STAT2} & 10 \\
\cline{3-5}
 &  & \multirow[t]{2}{*}{3} & {FOS $\wedge$ RELA $\wedge$ STAT2} & 9 \\
 &  &  & {RELA $\wedge$ STAT2 $\wedge$ STAT6} & 1 \\
\cline{1-5} \cline{2-5} \cline{3-5}
\multirow[t]{7}{*}{0.6} & \multirow[t]{7}{*}{1} & \multirow[t]{3}{*}{2} & {FOSL1 $\wedge$ STAT2} & 4 \\
 &  &  & {FOSL1 $\wedge$ RELA} & 3 \\
 &  &  & {RELA $\wedge$ STAT2} & 3 \\
\cline{3-5}
 &  & \multirow[t]{4}{*}{3} & {FOSL1 $\wedge$ RELA $\wedge$ STAT2} & 7 \\
 &  &  & {FOSL1 $\wedge$ IRF1 $\wedge$ RELA} & 1 \\
 &  &  & {FOS $\wedge$ FOSL1 $\wedge$ STAT2} & 1 \\
 &  &  & {FOS $\wedge$ RELA $\wedge$ STAT2} & 1 \\
\cline{1-5} \cline{2-5} \cline{3-5}
\bottomrule
\end{tabular}

%    \label{tab:threshold_cross_interpretable_mw_0.4_C_1}
%\end{table}
%
%\begin{table}[h]
%    \centering
%    \begin{tabular}{llllr}
\toprule
 &  &  &  & \textbf{Frequency} \\
$\mathbf{Q_y}$ & $\mathbf{C}$ & $\mathbf{D}$ & \textbf{TF combination} &  \\
\midrule
\multirow[t]{2}{*}{0.5} & \multirow[t]{2}{*}{1} & 2 & {RELA $\wedge$ STAT2} & 10 \\
\cline{3-5}
 &  & 3 & {FOS $\wedge$ RELA $\wedge$ STAT2} & 10 \\
\cline{1-5} \cline{2-5} \cline{3-5}
\multirow[t]{6}{*}{0.6} & \multirow[t]{6}{*}{1} & \multirow[t]{2}{*}{2} & {RELA $\wedge$ STAT2} & 5 \\
 &  &  & {FOSL1 $\wedge$ RELA} & 5 \\
\cline{3-5}
 &  & \multirow[t]{4}{*}{3} & {FOS $\wedge$ RELA $\wedge$ STAT2} & 4 \\
 &  &  & {FOSL1 $\wedge$ RELA $\wedge$ STAT2} & 3 \\
 &  &  & {FOSL1 $\wedge$ IRF1 $\wedge$ RELA} & 2 \\
 &  &  & {FOS $\wedge$ FOSL1 $\wedge$ RELA} & 1 \\
\cline{1-5} \cline{2-5} \cline{3-5}
\bottomrule
\end{tabular}

%    \label{tab:threshold_cross_interpretable_mw_0.6_C_1}
%\end{table}
%
%\begin{table}[h]
%    \centering
%    \begin{tabular}{llllr}
\toprule
 &  &  &  & \textbf{Frequency} \\
$\mathbf{Q_y}$ & $\mathbf{C}$ & $\mathbf{D}$ & \textbf{TF combination} &  \\
\midrule
\multirow[t]{2}{*}{0.5} & \multirow[t]{2}{*}{1} & 2 & {RELA $\wedge$ STAT2} & 10 \\
\cline{3-5}
 &  & 3 & {FOS $\wedge$ RELA $\wedge$ STAT2} & 10 \\
\cline{1-5} \cline{2-5} \cline{3-5}
\multirow[t]{6}{*}{0.6} & \multirow[t]{6}{*}{1} & \multirow[t]{2}{*}{2} & {RELA $\wedge$ STAT2} & 6 \\
 &  &  & {FOSL1 $\wedge$ RELA} & 4 \\
\cline{3-5}
 &  & \multirow[t]{4}{*}{3} & {FOS $\wedge$ RELA $\wedge$ STAT2} & 5 \\
 &  &  & {FOSL1 $\wedge$ IRF1 $\wedge$ RELA} & 2 \\
 &  &  & {FOS $\wedge$ FOSL1 $\wedge$ RELA} & 2 \\
 &  &  & {FOSL1 $\wedge$ RELA $\wedge$ STAT2} & 1 \\
\cline{1-5} \cline{2-5} \cline{3-5}
\bottomrule
\end{tabular}

%    \label{tab:threshold_cross_interpretable_mw_0.8_C_1}
%\end{table}
%
%\begin{table}[h]
%    \centering
%    \begin{tabular}{llllr}
\toprule
 &  &  &  & \textbf{Frequency} \\
$\mathbf{Q_y}$ & $\mathbf{C}$ & $\mathbf{D}$ & \textbf{TF combination} &  \\
\midrule
\multirow[t]{2}{*}{0.5} & \multirow[t]{2}{*}{1} & 2 & {RELA $\wedge$ STAT2} & 10 \\
\cline{3-5}
 &  & 3 & {FOS $\wedge$ RELA $\wedge$ STAT2} & 10 \\
\cline{1-5} \cline{2-5} \cline{3-5}
\multirow[t]{6}{*}{0.6} & \multirow[t]{6}{*}{1} & \multirow[t]{2}{*}{2} & {RELA $\wedge$ STAT2} & 6 \\
 &  &  & {FOSL1 $\wedge$ RELA} & 4 \\
\cline{3-5}
 &  & \multirow[t]{4}{*}{3} & {FOS $\wedge$ RELA $\wedge$ STAT2} & 5 \\
 &  &  & {FOSL1 $\wedge$ IRF1 $\wedge$ RELA} & 2 \\
 &  &  & {FOS $\wedge$ FOSL1 $\wedge$ RELA} & 2 \\
 &  &  & {FOSL1 $\wedge$ RELA $\wedge$ STAT2} & 1 \\
\cline{1-5} \cline{2-5} \cline{3-5}
\bottomrule
\end{tabular}

%    \label{tab:threshold_cross_interpretable_mw_1.0_C_1}
%\end{table}
%


\clearpage
\section{Greedy algorithm}\label{sec:threshold_greedy}
As an additional approach we explore a greedy algorithm with two goals in mind:
to investigate whether this simpler approach performs similarly,
and to verify that optimal rulesets for simple cases (low $C$ and $D$)
are similar to those produced by our earlier, more complicated approach.

In the greedy approach, we first find through brute force complete enumeration a rule with minimum \zo loss,
so with maximum train accuracy.
Then, the positively classified points are removed from the dataset and the cycle repeats.
Note that this greedy method is essentially a strategy to find a rule list rather than a rule set.
For $C=1$, the greedy approach finds an optimal ruleset.

Mathematically, we consider in step $\tau$ the points $\I_\tau$ and the corresponding positive points $\P_\tau$ and negative points $\N_\tau.$
Denote by rule $k_\tau$ the rule found in step $\tau$.
Then $\I_1 = \I$ and
\[\I_{\tau + 1} = \I_\tau \setminus \{i \in \I_\tau \fw k_\tau \in \K_{i}\} = \I_\tau \setminus \I_{\{k_\tau\}}, \tquad \tau \in \mathbb{N}_+.\]
Equivalently,
\[\I_{\tau + 1} = \{i \in \I_\tau \fw y_{\{k_\tau\}}(\bm{X}_i) \neq 1\}, \tquad \tau \in \mathbb{N}_+.\]

\begin{algorithm}[h!]
\caption{Greedy Ruleset Selection}
\label{alg:greedy}
\begin{algorithmic}[1]
\Require $\I$, $\bm{X}$, $\bm{y}$, $C$, $D$
\State $\I_1 \gets \I$
\For{$\tau = 1, \dots, C$}
    \State Find rule $k_\tau$ with at most $D$ clauses minimizing \zo loss w.r.t. $\I_\tau,$ $\bm{X},$ $y$
    \State $\I_{\tau+1} \gets \{i \in \I_{\tau} \fw \hat{y}_{\{k_\tau\}}(\bm{X}_i) \neq 1\}$
\EndFor \\
\Return $(k_\tau)_{\tau=1,\dots,C}$
\end{algorithmic}
\end{algorithm}

An advantage of the greedy algorithm is that it is easy to understand and to implement.
Furthermore, discretization via thresholding is done in every step separately.
This means that new, more relevant threshold values may be found in later steps,
when a part of the points is no longer considered.
A disadvantage is that there is no optimality guarantee for $C>1$.
Additionally, the Python implementation, while highly readable, is slow.
This is in part due to enumeration of all rules in each step, which is slow for larger values of $D$
(in our experiments, for $D\geq 3$).
Another aspect is the high-level nature of Python, hurting performance, despite the use of NumPy where possible.

The results for $D=1$ are summarized in \autoref{tab:greedy_yq_4_5_mw_0_C_3_D_1}.
While the greedy method was configured to continue after step two, there did not exist a rule that improved
performance.
The first rule provides reasonable performance.
The second rule provides barely any improvement.
For both $Q_y=0.5$ and $Q_y=0.6$, the same features appear in the rules.

The results for $D=2$ are summarized in \autoref{tab:greedy_yq_4_5_mw_0_C_3_D_2}.
In both tables, the value of adding another rule quickly diminishes.
The second and third rules barely increase accuracy.
The third rule found covers only a small fraction of the remaining positive points.
%Compared to \autoref{tab:threshold_interpretable_results_C_1} and \autoref{tab:threshold_interpretable_results_C_2},
%describing results of our earlier approach,
%the accuracy observed is similar.
Qualitatively, there is some overlap in the rules found, but there are differences.

The previous method (\autoref{sec:threshold-results}) performs similar to the greedy method.
Since for $C=1$ greedy is optimal, this suggests that the previous method is working as intended.
Sometimes, the greedy method outperforms the previous method.
There are two contributing factors for this.
First, the greedy method has more freedom in later steps.
Second, the previous approach is a heuristic, as it uses price-then-branch, and it was configured to use
only the heuristic pricer.
Nevertheless, overall performance is similar.

The conclusion of the experiments in this chapter is two-fold.
\begin{itemize}
    \item Compared to \autoref{sec:boolean-results}, decreasing complexity only moderately decreases performance.
    Instead, very simple rulesets can achieve a comparable accuracy, even when only using one clause per rule,
    which was not expected.
    It appears that the data is highly correlated and noisy, making it unsuitable for this off-the-shelf classifier model.
    \item The rulesets found are not consistent.
    The cross-validation experiment showed that rulesets vary among splits.
    The inconsistency of rules is supported by the wide variety of TF combinations observed in this chapter's tables.
\end{itemize}
We will therefore alter the model in an attempt to make it more robust and handle noise better, hopefully
providing more meaningful rulesets.

%\begin{table}
%    \centering
%    \bwtable{gen/PDL1ex2/tab_greedy_yq_4_mw_0_C_3_D_2}
%    \caption{}
%    \label{tab:greedy_yq_4_mw_0_C_3_D_2}
%\end{table}
%
%\begin{table}
%    \centering
%    \bwtable{gen/PDL1ex2/tab_greedy_yq_5_mw_0_C_3_D_2}
%    \caption{}
%    \label{tab:greedy_yq_5_mw_0_C_3_D_2}
%\end{table}

\begin{table}
    \centering
    \bwtable{gen/PDL1ex2/tab_greedy_yq_4_mw_0_C_3_D_1}
    \vspace{1em}
    \bwtable{gen/PDL1ex2/tab_greedy_yq_5_mw_0_C_3_D_1}
    \caption{Rules found in each step of the greedy algorithm, for $Q_y=0.5$ (top) and $Q_y=0.6$ (bottom).
    The full Broad dataset is used for training and $D=1$.
    In both cases, accuracy cannot be improved by any additional rule after step two.}
    \label{tab:greedy_yq_4_5_mw_0_C_3_D_1}
\end{table}

\begin{table}
    \centering
    \bwtable{gen/PDL1ex2/tab_greedy_yq_4_mw_0_C_3_D_2}
    \vspace{1em}
    \bwtable{gen/PDL1ex2/tab_greedy_yq_5_mw_0_C_3_D_2}
    \caption{Rules found in each step of the greedy algorithm, for $Q_y=0.5$ (top) and $Q_y=0.6$ (bottom).
    The full Broad dataset is used for training and $D=2$.}
    \label{tab:greedy_yq_4_5_mw_0_C_3_D_2}
\end{table}


%It can be formulated as follows:~
%\begin{subequations}
%    \label{eq:TP}
%    \begin{align}
%        \min \quad & \sum_{i \in \cN} \delta_i - \sum_{i \in \cP} \mu_i \delta_i +
%        \lambda \label{eq:TP-obj} \\
%        \st \quad & |J|\delta_i + \sum_{j \in \J} z_{ij} \leq |J|, & i \in \I, \label{eq:TP-con1}\\
%        & \delta_i +  \sum_{j \in J} z_{ij} \geq  1, & i \in \I,  \label{eq:TP-con2}\\
%        & z_{ij} \geq z_{i'j} & \forall\, i \in \I,\, j \in \J, \, :\, X_{ij} \leq X_{i'j}, \label{eq:TP-con3} \\
%        & \mathrlap{\delta\in \{0,1\}^\I,~z\in\{0,1\}^{\I\times \J}.}\label{eq:TP-vars}
%    \end{align}
%\end{subequations}
%Here $z_{ij}$ encodes whether point $i$ is cut off by feature $j$ in the chosen rule.
%Note that the rule complexity is not capped at $D$ in this formulation.
%In practice, this IP is hard to solve, due to its size.

